% New offline diagnostics. These are not part of the twenty-game LLM cohort.
\subsection{Reference-policy readiness and scenario validity}
\label{sec:agenda-reference-readiness}

The availability of a baseline implementation or a scenario library does not establish a validated comparator. Before using either to answer questions about strategic quality, adaptation or transfer, we audited the six supplied policies and all 130 static scenarios against the unchanged engine. We also conducted a small, separately registered local pilot of baseline-only games. These new executions are distinct from the historical twenty-game LLM cohort and contain no model queries.

\paragraph{Protocol and endpoints.}
We froze a local protocol before execution, using two purpose-chosen seeds, six homogeneous four-player rosters and four rotations of one mixed roster containing builder, always-buy, cash-conservative and random-legal policies. This yields twenty primary games. The engine remained uncapped; the evaluator stopped observation after 1,000 applied decisions and treated that endpoint as censoring, without assigning a net-worth winner. Invalid policy actions would stop a run and be reported, rather than trigger repair or fallback. We repeated two configurations with identical run identity and seed. The pilot tests execution, repertoire and censoring. Its small, selected seed set cannot estimate a population ranking or a seat effect. In particular, the random-legal policy derives its action randomization from the decision identity, including the run identifier; different rotation identifiers do not preserve identical policy draws.

\paragraph{Live-policy results.}
The twenty games applied 19,494 decisions and emitted 179,816 events with no invalid generated action and no fallback. Nineteen games reached the observation limit. The remaining mixed-roster game terminated after 494 decisions with a builder-policy win. No policy initiated a trade in the pilot, and only 44 construction actions occurred. The two same-configuration repeats matched complete event streams, action streams and final snapshots. A separate replay of all twenty retained action streams reproduced every event and final snapshot. These checks establish reproducibility under this implementation, not strategic competence or an absence of eventual termination in the censored games.

\paragraph{Static fixtures expose different failures.}
For each supplied scenario, we constructed each policy's action and checked its economic legality at the supplied decision and board state, using the same engine validator used by the fixture runner. Of 780 policy--scenario choices, 75 were rejected. The six policies' counts are shown in Table~\ref{tab:agenda-fixture-legality}. Eleven of the 130 supplied reference actions were also rejected. Examples include construction on an incomplete or split-owned red group, a purchase or auction for an already-owned deed, a hotel downgrade with only two bank houses, and a proposed transfer involving a deed owned by a third player. These checks validate the action against the fixture's current pending decision; they do not reconstruct a historical path or hidden random state.

\begin{table}[htbp]
\centering
\caption{Economic legality of generated actions on the 130 supplied static scenarios. These are fixture-conditioned results, not live-game failure rates. The pilot's live decisions had no invalid actions.}
\label{tab:agenda-fixture-legality}
\begin{tabular}{lrr}
\toprule
Policy & Legal & Rejected \\
\midrule
Always buy & 114 & 16 \\
Auction aggressive & 122 & 8 \\
Builder & 107 & 23 \\
Cash conservative & 122 & 8 \\
No trade & 122 & 8 \\
Random legal & 118 & 12 \\
\midrule
Total & 705 & 75 \\
\bottomrule
\end{tabular}
\end{table}

The rejection counts should not be attributed wholesale to baseline strategy. Many fixture option lists name properties that the supplied board does not permit the actor to mortgage or develop. The policies normally rely on engine-generated options, and the live pilot produced no such invalid actions. A follow-up static-state diagnostic flagged 55 scenarios for at least one inconsistency under the standard stock-conservation or development conditions. All 55 had a house-stock discrepancy; overlapping flags included split ownership, mortgages within developed groups, a hotel-stock discrepancy, and ten liquidation shortfalls inconsistent with stated cash and debt. A deliberately nonconserving scarcity intervention might be a useful artificial task, but its altered condition must be declared. Sixty-eight scenarios pass both the reference-action and specified static-invariant screens. This identifies candidates for independent task adjudication, not certified capability probes: passing this small checklist does not prove state reachability or reference-action optimality.

The scoring library also measures more than executed strategy: 128 of 130 scenarios assign points to private-note keyword matches, with a median weight of 15.8\% and a maximum of 25.7\%. Such criteria can measure expression of prescribed concepts; they are not independent evidence of sound strategic choices. The supplied preferred actions remain hand-authored preferences rather than certified expected-utility optima. Forty-one scenarios induce identical action names and arguments from all six policies, further limiting their role as diverse comparators. In forced proposal fixtures the policies can emit empty bundles; none constructs a meaningful bargaining strategy.

These diagnostics answer a concrete validity question: the current library and policies cannot simply be relabeled as a validated strategic-capability battery. They supply executable infrastructure, but further scenario adjudication, opponent design and target definition are required. No reference fixture or production policy was repaired during this audit. The engine source is byte-identical to the engine bundled in the historical executable audit, so the reported fixture discrepancies do not arise from using a different engine version.

\subsection{A design for the unanswered experimental questions}
\label{sec:agenda-prospective-evaluation}

The remaining questions require distinct targets, rather than one aggregate competence score. The following arms are specified in a machine-readable protocol and have \emph{not} been run. The offline runner has no provider adapter, rejects nonzero-spend or network-enabled configurations, and blocks socket connections. The purpose is to make the proposed research inspectable before any new model expenditure or task-condition change.

\paragraph{Local strategic execution and economic value.}
A local task should first declare its endpoint, such as paying a current debt or reaching a specified construction allocation. Related states and source games should be assigned together to development or held-out sets. Engine witnesses can certify feasibility; they cannot certify that the target maximizes eventual winning probability. An action-value study additionally needs a utility, horizon and opponent distribution, with repeated stochastic continuations. We would compare endpoint attainment, encoding failures and execution separately from independently scored explanation content. General regret or exploitability claims require a defined comparator and opponent class. The existing historical states are public and heavily inspected, so evaluation on them alone cannot establish unfamiliar-task transfer.

\paragraph{Adaptation and memory.}
An adaptation experiment should randomize relevant feedback against a matched control and then present comparable later opportunities, followed by a held-out opponent-policy shift. The unit is a complete task sequence, not an individual correction. The treatment effect is improvement after feedback relative to control, with state-driven reactions, invalid outputs and missing opportunities retained. Cross-game adaptation requires an explicit memory-carryover condition. Weight learning requires training. Neither follows from repeated games or successful corrective retries. Feedback and memory manipulations change the model-facing condition and therefore require separate, versioned authorization; they are not changes to the historical prompt.

\paragraph{Opponent transfer and chance.}
Cross-play should vary opponent families and independently sampled seeds while rotating seats within seed blocks. Opponents must first pass the quality checks above and include meaningful negotiation and development behavior. Complete-system performance and a base-model comparison are different estimands; response policies should be controlled or deliberately factorialized. A historical seed does not guarantee matched opportunities once actions alter random-draw consumption. We would report wins, censoring, rule/encoding failures and task-level endpoints, then assess out-of-sample ordering on held-out opponents. Two chosen pilot seeds do not determine a sufficient sample size or justify a confidence interval for the LLM cohort.

\paragraph{General capability and multi-agent capability.}
Question B requires independent measurements of general ability across multiple tiers within several families. Model configuration, inference resources and response policy must be recorded or controlled. The outcome is a prespecified multi-agent battery, distinct from honesty or instruction adherence. A predictive relationship should be tested on held-out models or families and compared with a family/resource-only baseline. Adding four external leaderboard scores to the existing four configurations would not identify the requested relationship. Precision and expected effect size, rather than an arbitrary game-count target, must determine the prospective sampling plan.

\paragraph{Framing and coordination.}
An optional message experiment can randomize truthful complete disclosure, omission and neutral wording while holding economic terms and other delivered facts fixed. Factual endorsement, action choice and subsequent payoff are separate outcomes. Treatment effects would concern the specified manipulation, not a latent intent label. Changes to public messages are new task conditions, not retroactive changes to historical interactions. A communication intervention can test a specified collective mechanism; parallel actions under shared incentives are not by themselves proof that communication caused a coalition.

\paragraph{Cost and execution gate.}
The saved historical cost table records \$141.71 of known OpenRouter charges and 1,603 attempts with missing cost, using rounded source aggregates. Those values describe an earlier workload and are not current prices or safe upper bounds. A paid campaign must first freeze its roster, task conditions, outcome definitions, stopping rule and maximum expenditure. Before each request, a separate runner would need to reserve a defensible upper cost from current prices and billable-token limits, including retries, and refuse any request whose bound is unknown or exceeds the remaining allowance. No such provider execution is enabled in this package. Until these designs are executed and validated, learning, unfamiliar-opponent generalization and a general-to-multi-agent capability relationship remain empirical gaps, not completed results.
